\begin{abstract}
La convergencia de computación cuántica y constelaciones de satélites LEO/MEO plantea amenazas existenciales a la seguridad de las comunicaciones espaciales basadas en criptografía de clave pública clásica. Este artículo propone un marco integral para una infraestructura orbital con resiliencia cuántica mediante la integración sistemática de algoritmos de criptografía poscuántica (PQC) estandarizados por NIST (ML-KEM, ML-DSA, SLH-DSA) en protocolos de redes satelitales de próxima generación. Se analizan restricciones de recursos (energía, ancho de banda, radiación) en entornos espaciales y se presenta una arquitectura híbrida PQC-clásica con agilidad criptográfica. Mediante modelado analítico, simulaciones y evaluación en plataformas emuladas de satélites, se demuestran mejoras en resiliencia contra ataques harvest-now-decrypt-later mientras se mantiene overhead aceptable en latencia y throughput. Los resultados indican viabilidad para despliegues en constelaciones mega-LEO, contribuyendo a estándares CCSDS y arquitecturas 6G-NTN seguras. Se discuten desafíos de implementación, verificación formal y hoja de ruta de migración 2025-2035.
\end{abstract}

\textbf{Keywords:} post-quantum cryptography, PQC, satellite networks, LEO/MEO constellations, quantum resilience, space security, orbital infrastructure, cryptographic agility, satellite protocols, 6G-NTN.